\newpage % Nova página
\thispagestyle{empty} % Remove o número da página
\clearpage % Finaliza a página atual

\begin{center}
    \section*{Abstract} % Centraliza o título "Abstract" (sem numeração)
\end{center}

\par Supervisory control and data acquisition (SCADA) systems are the software layer responsible for monitoring, supervising and controlling industrial processes in real time, bringing together in a single interface the reading of field variables, historical logging, alarm signaling and the dispatch of commands to the equipment. Present in sectors such as water and sanitation, energy, manufacturing and oil and gas, these systems allow the operator to follow the plant remotely, identify anomalies and act on the process without physical travel, which reduces operating costs and increases production reliability. Conventional SCADA systems, however, require high investment in proprietary licenses, dedicated hardware and local infrastructure, as well as specialized installation and maintenance, which keeps them out of reach for small and medium-sized enterprises. In view of this, this work presents the development of a service-oriented web supervisory system, running on cloud infrastructure. The proposal prioritizes usability. Access is provided through the browser, with no client-side installation, a single infrastructure serves multiple plants, and access to screens and actions is controlled by role, following the principle of least privilege. Within the context of Industry 4.0, this digitalization of the supervision and management layers delivers as a service functionalities once restricted to the physical environment of the plant. The solution adopts a microservices architecture, with a Next.js web application that concentrates the supervision interface and the backend for registry, history, alarms and commands, a scheduler that coordinates the periodic readings, and an acquisition service that communicates with the field equipment through the Modbus TCP/RTU protocol. History is persisted in a relational database (PostgreSQL), the current value of each point is kept in an in-memory cache (Redis), and the services communicate through a message bus (Moleculer), hosted on managed cloud services (Amazon ECS, RDS and ElastiCache) to obtain scalability and redundancy. The concepts of Internet of Things (IoT), Industrial Internet of Things (IIoT), cyber-physical systems (CPS) and Industry 4.0 that underpin the proposal are reviewed, and the classic SCADA requirements, data acquisition and logging, time synchronization, access levels, command and control, scalability and redundancy, are mapped onto the implemented components. Ultimately, the monitoring module is expected to be validated through tests and a demonstrative scenario, evidencing the correct access to process variables and the stability of the solution under the expected data volume.

\leavevmode\newline
\par\textbf{Keywords:} SCADA; industrial process supervision; remote monitoring; cloud computing; service-oriented architecture; Industry 4.0; Modbus. % Negrito para "Keywords", seguido das palavras-chave em inglês

\thispagestyle{empty} % Remove a numeração da página
\clearpage % Nova página após o abstract
